% Compile with make report-pdf from the repository root, or:
%   python build_latex.py report
% Fill every TODO and replace all -- cells with measured results.
% Maximum: 5 pages INCLUDING figures and first-page team information;
% references are excluded. Do not change the CVPR font sizes or margins.
\documentclass[10pt,twocolumn,letterpaper]{article}
\usepackage[pagenumbers]{cvpr}
\usepackage[T1]{fontenc}
\usepackage{microtype}
\usepackage[hidelinks]{hyperref}
\graphicspath{{figures/}}
\newcommand{\todo}[1]{\textcolor{red}{[TODO: #1]}}
\newcommand{\missing}{--}

\title{Semi-Supervised Stopping-Reason Prediction\\and Clinical-Trial Outcome Classification}
\author{\todo{Team member 1} \qquad \todo{Team member 2} \qquad \todo{Team member 3}\\
CS5785 / ORIE5750 / ECE5414 --- Midterm Project, Fall 2026}

\begin{document}
\maketitle
% The assignment asks for team responsibilities on the cover page. Keep these
% on the first page with the title; a separate cover would consume page budget.
\noindent\textbf{Team responsibilities.}
\todo{Member 1: contribution; member 2: contribution; member 3: contribution.}

\begin{abstract}
\todo{Briefly describe both tasks, the approaches actually tested, the best
validation results, and the main insight. Write this after completing the report.}
\end{abstract}

\section{Methods}
\subsection{Data and preprocessing}
\todo{Describe the provided splits, missing labels (-100), class imbalance,
feature choices, scaling of reason embeddings, missing drug data, and
preprocessing. Fit learned transformations on training data and describe
validation-based selection. Use IDs only for matching.}

\subsection{Part 1: Semi-supervised reason labeling}
\todo{Describe at least one unsupervised method and how it generates the five
missing reason labels. Explain initialization, cluster count or neighborhood
parameters, convergence criteria, iteration limits, random seeds, and how
clusters map to multiple labels. Describe any supervised stage, thresholds,
feature ablations, and every method tested. Use no deep learning.}

\subsection{Part 2: Outcome classification}
\todo{Describe the original classifier, phase handling, features, tuning,
probabilities, and F1 thresholds. Reason embeddings and reason labels must
be used only for error analysis, never as predictive inputs to this classifier.}

\subsection{Error analysis and modification}
\todo{Match unsuccessful validation trials to Part 1 using trial IDs. Distinguish
true reason labels from Part 1 predictions; treat unmatched trials as a
separate no-recorded-stopping-reason group. Count each trial once per phase;
multi-label trials may belong to several reason groups. Explain the resulting
classifier modification and its motivation.}

\section{Results}
\subsection{Part 1 validation performance}
\todo{Compare all tested approaches and feature choices on part1\_val.csv.
Report sample-wise F1, macro F1, and F1 for each of the five labels.}
\begin{table}[ht]
\centering
\begin{tabular}{lcc}
\toprule
Metric & Baseline & Selected \\
\midrule
Sample-wise F1 & \missing & \missing \\
Macro F1 & \missing & \missing \\
Safety F1 & \missing & \missing \\
Funding F1 & \missing & \missing \\
Efficacy F1 & \missing & \missing \\
Recruitment F1 & \missing & \missing \\
Other F1 & \missing & \missing \\
\bottomrule
\end{tabular}
\caption{Part 1 validation results. Add comparisons for other tested models
and feature ablations; identify the models in the column headings.}
\label{tab:part1}
\end{table}

\subsection{Part 2 validation performance}
\todo{Report metrics for each phase, treating success (label=1) as positive.
Compute F1 from binary predictions and ROC-AUC and average precision from
probabilities. Compare the original and modified classifiers even if F1 falls.}
\begin{table}[ht]
\centering
\begin{tabular}{llccc}
\toprule
Phase & Model & F1 & ROC-AUC & AP \\
\midrule
I & Original & \missing & \missing & \missing \\
I & Modified & \missing & \missing & \missing \\
II & Original & \missing & \missing & \missing \\
II & Modified & \missing & \missing & \missing \\
III & Original & \missing & \missing & \missing \\
III & Modified & \missing & \missing & \missing \\
\bottomrule
\end{tabular}
\caption{Part 2 validation performance. AP is PR-AUC measured as average precision.}
\label{tab:part2}
\end{table}

\subsection{Errors by stopping reason}
\todo{For every phase and reason group, report the number of unsuccessful
validation trials and the fraction incorrectly predicted as successful, for
both classifiers. State label provenance and explain small or empty groups;
use N/A for undefined rates.}
% Duplicate this table for phases II and III, or combine all phases in a table*.
\begin{table}[ht]
\centering
\begin{tabular}{lrrr}
\toprule
Reason (Phase I) & Count & Orig. rate & Mod. rate \\
\midrule
Safety & \missing & \missing & \missing \\
Funding & \missing & \missing & \missing \\
Efficacy & \missing & \missing & \missing \\
Recruitment & \missing & \missing & \missing \\
Other & \missing & \missing & \missing \\
No recorded reason & \missing & \missing & \missing \\
\bottomrule
\end{tabular}
\caption{False-success fractions among label=0 trials. Repeat for phases II
and III. Groups overlap when trials have multiple reasons.}
\label{tab:reason-errors}
\end{table}

% Example figure: export it from a submitted notebook into figures/.
% \begin{figure}[ht]
%   \centering
%   \includegraphics[width=\linewidth]{feature_ablation.pdf}
%   \caption{Describe the comparison, split, axes, and takeaway.}
%   \label{fig:ablation}
% \end{figure}

\section{Discussion}
\todo{Explain why particular methods, hyperparameters, and features helped or
failed, including why including every feature may hurt. Discuss temporal
shift and the validation/test gap (Part 2 Kaggle uses probability log loss),
reason-category sizes, possible errors in predicted reason labels, and label
noise. Evaluate the modification's trade-offs, limitations, and main insights.}

% Cite sources where used with \cite{key}; add corresponding \bibitem entries.
% Include all papers, libraries, websites, and public code used for ideas,
% code, or phrasing. Do not submit the placeholder reference below.
\begin{thebibliography}{9}
\bibitem{replace-me}
\todo{Replace with actual references: author, title, venue/year or URL and
access date, as appropriate.}
\end{thebibliography}
\end{document}
